\begin{exo}{}
	Le plan est muni d'un repère orthonormé $\vOIJ$. On considère les points suivants :
	\vspace*{-2mm}
	$$A\coordl{-3}{1}\;,\;B\coordl{1}{3}\;\text{et}\;C\coordl{4}{-1}$$

	\begin{enumerate}\vspace*{-3mm}
		\item Dans un repère, placer les points $A$, $B$ et $C$.
		\item Calculer les coordonnées de $\vec{AB}$.
		\item Placer le point $D$ tel que le vecteur $\vec{DC}$ ait pour coordonnées $\coord{4}{2}$. Lire graphiquement les coordonnées du point $D$.
		\item Démontrer par le calcul que $\vec{AB}=\vec{DC}$.
		\item En déduire la nature exacte du quadrilatère $ABCD$. Justifier la réponse.
	\end{enumerate}
\end{exo}
